\subsection{Inductive Voltage Divider}

This block is constructed from two inductors in series and is used
to attenuate a voltage signal.

\subsubsection{Schematic}
\begin{circuitfig}
	% Paths, nodes and wires:
	\node[ground, rotate=-90] at (5.5, 5.25){};
	\node[ground, rotate=-90] at (5.5, 4.75){};
	\draw (5.5, 6.75) to[sinusoidal voltage source, l={$V_{cc}$}] (5.5, 5.25);
	\draw (5.5, 4.75) to[sinusoidal voltage source, l={$V_{ee}$}] (5.5, 3.25);
	\draw (7, 5) -- (7.75, 5);
	\draw (7, 7) -| (5.5, 7) -| (5.5, 6.75);
	\node[shape=rectangle, minimum width=0.715cm, minimum height=0.465cm] (N1) at (8.125, 5){} node[anchor=center] at (N1.text){$V_{out}$};
	\draw (7, 5.25) -| (7, 5);
	\draw (7, 4.75) -| (7, 5);
	\draw (5.5, 3.25) -| (5.5, 3) -- (7, 3);
	\draw (7, 7) to[american inductor, l={$L_1$}] (7, 5.25);
	\draw (7, 4.75) to[american inductor, l={$L_2$}] (7, 3);
\end{circuitfig}

\subsubsection{Transfer Function}
\textbf{Small Signal, Frequency Domain, AC Input Analysis (Q4)}: KCL at
$V_{out}$ (in the \textbf{Large Signal, Time Domain, DC Input Analysis (Q5)}
both inductors are shorts and the divider shorts $V_{cc}$ to $V_{ee}$):
\[
  \frac{V_{out} - V_{cc}}{sL_1} + \frac{V_{out} - V_{ee}}{sL_2} = 0
\]

\paragraph{Superposition:}
\begin{enumerate}
  \item Due to \(V_{cc}\) alone (\(V_{ee} = 0V\)):
    \begin{align*}
      \frac{V_{out} - V_{cc}}{sL_1} + \frac{V_{out} - 0}{sL_2} &= 0 \\
      \frac{V_{out} - V_{cc}}{sL_1} + \frac{V_{out}}{sL_2} &= 0 \\
      \frac{V_{out}}{sL_1} - \frac{V_{cc}}{sL_1} + \frac{V_{out}}{sL_2} &= 0 \\
      V_{out}\left(\frac{1}{sL_1} + \frac{1}{sL_2}\right) - \frac{V_{cc}}{sL_1} &= 0 \\
      V_{out}\left(\frac{L_2 + L_1}{sL_1 L_2}\right) - \frac{V_{cc}}{sL_1} &= 0 \\
      V_{out}\left(\frac{L_2 + L_1}{sL_1 L_2}\right) &= \frac{V_{cc}}{sL_1} \\
      sL_1 L_2 \cdot V_{out}\left(\frac{L_2 + L_1}{sL_1 L_2}\right) &= sL_1 L_2 \cdot \frac{V_{cc}}{sL_1} \hspace{1cm} \text{(Multiplying both sides by \(sL_1 L_2\))} \\
      V_{out}(L_2 + L_1) &= V_{cc} L_2 \\
      V_{out} &= V_{cc} \frac{L_2}{L_1 + L_2} \\
      H_1(s) = \frac{N_1(s)}{D_1(s)} = \frac{V_{out_1}}{V_{cc}} &= \frac{L_2}{L_1 + L_2}
    \end{align*}
  \item Due to \(V_{ee}\) alone (\(V_{cc} = 0V\)):
    \begin{align*}
      \frac{V_{out} - 0}{sL_1} + \frac{V_{out} - V_{ee}}{sL_2} &= 0 \\
      \frac{V_{out}}{sL_1} + \frac{V_{out} - V_{ee}}{sL_2} &= 0 \\
      \frac{V_{out}}{sL_1} + \frac{V_{out}}{sL_2} - \frac{V_{ee}}{sL_2} &= 0 \\
      V_{out}\left(\frac{1}{sL_1} + \frac{1}{sL_2}\right) - \frac{V_{ee}}{sL_2} &= 0 \\
      V_{out}\left(\frac{L_2 + L_1}{sL_1 L_2}\right) - \frac{V_{ee}}{sL_2} &= 0 \\
      V_{out}\left(\frac{L_2 + L_1}{sL_1 L_2}\right) &= \frac{V_{ee}}{sL_2} \\
      sL_1 L_2 \cdot V_{out}\left(\frac{L_2 + L_1}{sL_1 L_2}\right) &= sL_1 L_2 \cdot \frac{V_{ee}}{sL_2} \hspace{1cm} \text{(Multiplying both sides by \(sL_1 L_2\))} \\
      V_{out}(L_2 + L_1) &= V_{ee} L_1 \\
      V_{out} &= V_{ee} \frac{L_1}{L_1 + L_2} \\
      H_2(s) = \frac{N_2(s)}{D_2(s)} = \frac{V_{out_2}}{V_{ee}} &= \frac{L_1}{L_1 + L_2}
    \end{align*}
\end{enumerate}
\paragraph{Total Output:}
\begin{align*}
  V_{out} &= V_{out_1} + V_{out_2} \\
  V_{out} &= V_{cc} \frac{L_2}{L_1 + L_2} + V_{ee} \frac{L_1}{L_1 + L_2}
\end{align*}

\subsubsection{Analysis}
\begin{enumerate}
  \item \textbf{Inductor Ratio (\(\frac{L_2}{L_1}\)):}
    From KCL at \(V_{out}\):
    \begin{align*}
      \frac{V_{out} - V_{cc}}{sL_1} + \frac{V_{out} - V_{ee}}{sL_2} &= 0 \\
      \frac{V_{out} - V_{cc}}{sL_1} - \frac{V_{ee} - V_{out}}{sL_2} &= 0 \hspace{1cm} \text{(Writing \(V_{out} - V_{ee} = -(V_{ee} - V_{out})\))} \\
      \frac{V_{out} - V_{cc}}{sL_1} &= \frac{V_{ee} - V_{out}}{sL_2} \\
      (V_{out} - V_{cc}) sL_2 &= (V_{ee} - V_{out}) sL_1 \hspace{1cm} \text{(Multiplying both sides by \(sL_1 \cdot sL_2\))} \\
      \frac{sL_2}{sL_1} &= \frac{V_{ee} - V_{out}}{V_{out} - V_{cc}} \hspace{1cm} \text{(Dividing both sides by \(sL_1(V_{out} - V_{cc})\))} \\
      \frac{L_2}{L_1} &= \frac{V_{ee} - V_{out}}{V_{out} - V_{cc}}
    \end{align*}
  \item \textbf{Impedances:}
    \paragraph{At Input:} Looking in from \(V_{cc}\) with \(V_{ee}\) shorted,
    \(L_1\) and \(L_2\) are in series:
      \begin{align*}
        Z_{in} &= sL_1 + sL_2 \\
        Z_{in} &= s(L_1 + L_2)
      \end{align*}
    \paragraph{At Output:} Looking in from \(V_{out}\) with both sources
    shorted, \(L_1\) and \(L_2\) are in parallel:
      \begin{align*}
        Z_{out} &= \frac{1}{\frac{1}{sL_1} + \frac{1}{sL_2}} \\
        Z_{out} &= \frac{1}{\frac{L_2 + L_1}{sL_1 L_2}} \\
        Z_{out} &= \frac{sL_1 L_2}{L_1 + L_2}
      \end{align*}
  \item \textbf{Loading Effect:} \\
    Similar to resistive dividers.
\end{enumerate}

\subsubsection{Question}
Design an inductive voltage divider that attenuates a 10V signal to 7V.
Establish the ratio of the inductors and calculate the input and output
impedance. If one of the inductors is 18 nH, calculate the other inductor's
value.
\paragraph{Solution:}
\begin{enumerate}
  \item Calculate the inductor ratio with \(V_{cc} = 10V\), \(V_{out} = 7V\)
    and \(V_{ee} = 0V\):
    \begin{align*}
      \frac{L_2}{L_1} &= \frac{V_{ee} - V_{out}}{V_{out} - V_{cc}} \\
      \frac{L_2}{L_1} &= \frac{0V - 7V}{7V - 10V} \\
      \frac{L_2}{L_1} &= \frac{-7V}{-3V} \\
      \frac{L_2}{L_1} &= \frac{7}{3} \approx 2.33
    \end{align*}
  \item Calculate the input impedance:
    \begin{align*}
      Z_{in} &= s(L_1 + L_2) \\
      Z_{in} &= s\left(L_1 + \frac{7}{3}L_1\right) \hspace{1cm} \text{[substituting \(L_2 = \frac{7}{3}L_1\)]} \\
      Z_{in} &= sL_1\left(\frac{3}{3} + \frac{7}{3}\right) \\
      Z_{in} &= s\left(\frac{10}{3}L_1\right) \approx s(3.33L_1)
    \end{align*}
  \item Calculate the output impedance:
    \begin{align*}
      Z_{out} &= \frac{sL_1 L_2}{L_1+L_2} \\
      Z_{out} &= \frac{sL_1\left(\frac{7}{3}L_1\right)}{L_1 + \frac{7}{3}L_1} \hspace{1cm} \text{[substituting \(L_2 = \frac{7}{3}L_1\)]} \\
      Z_{out} &= \frac{s\left(\frac{7}{3}L_1^2\right)}{\frac{10}{3}L_1} \\
      Z_{out} &= s\left(\frac{7}{3} \cdot \frac{3}{10}\right)L_1 \\
      Z_{out} &= s(0.7L_1)
    \end{align*}
  \item Calculate the other inductor's value (taking \(L_1 = 18nH\)):
    \begin{align*}
      L_2 &= \frac{7}{3}L_1 \\
      L_2 &= \frac{7}{3}(18nH) \\
      L_2 &= 42nH
    \end{align*}
    so that
    \begin{align*}
      Z_{in} &= s\left(\frac{10}{3} \cdot 18nH\right) = s(60nH) = j 2 \pi f (60nH) \\
      Z_{out} &= s(0.7 \cdot 18nH) = s(12.6nH) = j 2 \pi f (12.6nH)
    \end{align*}
\end{enumerate}
